% Notes of the Montreal seminar, 1970 (cote 129, typed part): archive cote 129, site page 38 (PDF page 39).
% typescript. Draft compared: none.
% AI-transcribed by Claude Opus 5.5 (claude-opus-5-5, adaptive thinking, effort high), Message Batches API, batch mtl129_full (msgbatch_01AGnDKHwy6ATCM7ehUa6263), 2026-11-13; not yet checked by eye.
\scanpage{129/38}
\mspage{\hand{II.10}}

On en déduit donc que, si $G$ est localement de type fini sur
$S$ , \unsure{$\mathcal{F}_{G/S}$|$\hat{\mathcal{F}}_{G/S}$} est localement nilpotent si et seulement si $\mathcal{J}$ est localement
nilpotent. Pour un groupe fini localement libre sur $S$ cela équivaut au
fait qu'il est infinitésimal.

\uline{Démonstration de 2.5.4.}

On montre d'une manière générale (SGA 5, XIV (Houzel)) que si
$G$ est un schéma localement de présentation finie sur $S$ , \unsure{$\mathcal{F}_{G/S}$|$\hat{\mathcal{F}}_{G/S}$} est un
isomorphisme si et seulement si $G$ est étale sur $S$ .

3. \uline{Décomposition d'un p-groupe fini localement libre sur une base de carac-}
\uline{téristique} $p$ \uline{réduite à un point} $s$ .

3.1. $G$ étant un groupe fini localement libre sur $S$ , on a une sui-%
te exacte canonique, fonctorielle en $G$
\[ (*) \qquad 0 \rightarrow G^{0} \rightarrow G \rightarrow G^{et} \rightarrow 0 \]
où $G^{0}$ est la composante neutre de $G$ et $G^{et} = G/G^{0}$ . Le sous-groupe $G^{0}$
est ouvert dans $G$ (SGA$_{3}$, VI$_{B}$, 3.9), il est donc fini localement libre sur
$S$ et c'est le plus grand sous-groupe infinitésimal de $G$ . Le groupe quo-%
tient $G^{et}$ est étale, (SGA 3, VI$_{A}$, 5.5) et il est fini localement libre
sur $S$ , d'après 1.6.

Appliquant la même construction à $G^{*}$ , on obtient un sous-grou-%
pe $G^{mult} = (G^{*}/(G^{*})^{0})^{*}$ , qui est le plus grand sous-groupe de type multi-%
plicatif de $G$ .

\mspage{\hand{38}}
